% Abstract only. Count visible headings, text, spaces and paragraph breaks.
\textbf{Background.} Breathing and visual feedback may influence conscious experience and the perceived boundary between body and environment.

\textbf{Aims.} Develop breath-responsive virtual reality (VR) and assess changes in subjective experience and peripersonal space.

\textbf{Method.} Viscereality links breathing to visual changes. A controlled study tested 47 participants; 39 were included after breathing-signal exclusions. Participants completed two VR conditions and a minimal-feedback control during box breathing. A separate audio-tactile study examined breathing phase and responses to approaching sounds; 40 of 42 released records entered its primary behavioural analysis. A later affective-feedback study remains incomplete.

\textbf{Results.} Compared with control, VR improved synchronization with the breathing guide and increased reported spatial self-extension and several positive and perceptual experiences; distress also increased modestly. The two VR conditions did not differ reliably. Behavioural reaction-time patterns did not establish a reliable change in the peripersonal-space boundary. The later study has eight participants with both questionnaire mappings and three partial records, against a target of 24. Outputs include two conference papers, a preprint and PPS Kit software.

\textbf{Conclusions.} Breath-responsive VR can alter reported experience. The findings do not establish clinical benefit or the proposed spatial mechanism. Research remains unfinished; journal acceptance is not established.
